% Generated draft: requires mathematical and source review.
\documentclass[11pt]{article}
\usepackage[margin=25mm]{geometry}
\usepackage{fontspec}
\setmainfont{Latin Modern Roman}
\usepackage{amsmath,amssymb,mathtools,unicode-math}
\setmathfont{Latin Modern Math}
\usepackage{enumitem,needspace,xurl,hyperref}
\setlength{\parindent}{0pt}
\setlength{\parskip}{5pt}
\setlength{\emergencystretch}{3em}
\newcommand{\N}{\mathbb{N}}
\newcommand{\Z}{\mathbb{Z}}
\newcommand{\Q}{\mathbb{Q}}
\newcommand{\R}{\mathbb{R}}
\newcommand{\C}{\mathbb{C}}
\newcommand{\F}{\mathbb{F}}
\newcommand{\A}{\mathbb{A}}
\newcommand{\PP}{\mathbb{P}}
\newcommand{\E}{\mathbb{E}}
\newcommand{\eps}{\varepsilon}
\newcommand{\dd}{\,\mathrm{d}}
\begin{document}
\section*{UnsolvedMath ID: 6200065}
\textbf{Source problem code:} catalogue entry
\subsection{Definitions and mathematical statement}
For a discrete Kleinian group \(\Gamma\le\operatorname{PSL}_2(\mathbb C)\), acting on hyperbolic three-space \(\mathbb H^3\), its limit set \(\Lambda(\Gamma)\subseteq\partial_\infty\mathbb H^3\cong S^2\) is the set of boundary accumulation points of an orbit \(\Gamma x\), \(x\in\mathbb H^3\). A space \(L\) is locally connected if each point and each neighborhood of that point contain an open connected neighborhood of the point, relative to \(L\).

The supplied statement begins with the assertion
\[
 \Lambda(\iota(G))\text{ is locally connected},
\]
but neither \(G\) nor the representation \(\iota\) or its hypotheses are supplied. It then says that in the presence of ``two geodesic laminations'' the limit set is the whole sphere,
\[
 \Lambda(\iota(G))=S^2,
\]
and that local connectedness is consequently not a useful discriminator in that situation. These are contextual source assertions, not universal claims for arbitrary Kleinian groups: the lamination setting and its relation to the representation are absent.

Most importantly, the excerpt ends with ``Then the correct reformulation is as follows:'' and contains no reformulation after the colon. The intended problem, its quantifiers, and any nontrivial replacement for local connectedness are therefore missing. No candidate reformulation or restriction to a guessed surface-group setting is inserted.

The research summary connects the topic to the Cannon--Thurston/local-connectivity program and mentions Kapovich--Kleiner conditions, but provides no precise group class, theorem, or citation that restores the missing question. Its general open/partial assessment and the 2005 problem label are retained as source information, not an independently verified current status or a theorem covering the displayed assertions.
\subsection{Short English statement}
The record discusses local connectedness of a Kleinian limit set and says a different formulation is needed when the limit set is the whole sphere. That reformulation is missing, along with the group, representation, and lamination hypotheses, so the intended question cannot be recovered from this excerpt.
\subsection{Sources}
\begin{itemize}
\item Local registry: \url{data/unsolvedmath/problems.json}, permanent UnsolvedMath ID \(6200065\), source code catalogue entry; complete supplied statement, background, and research summary, updated 2026-08-06.
\item Kapovich, \emph{Problems on Boundaries of Groups and Kleinian Groups} (2005), Problem 65, PDF page 18: \url{https://www.math.ucdavis.edu/~kapovich/EPR/problems.pdf}. Supplied normalized PDF-text extraction accessed 2026-07-29; current status marked as needing review.
\item The summary names the Cannon--Thurston program and Kapovich--Kleiner work without full publications or exact hypotheses. No source for the missing reformulation is separately supplied.
\end{itemize}
\end{document}
